\documentclass[11pt]{article}
\usepackage{mathunal-base}
\mathunalfuente{serif}

\renewcommand{\examtitulo}{Segundo Examen Parcial de Álgebra Lineal}
\renewcommand{\examfecha}{29 de octubre de 2022}

\newcommand{\vf}{\fbox{V}\ \fbox{F}}
\newcommand{\sinosino}{\fbox{$SI$}\hspace{0.5cm}\fbox{$NO$}}

\begin{document}

\examheader

\vspace{4pt}
\noindent
\begin{minipage}[t]{0.62\linewidth}
\textbf{Puntaje. Sólo para uso Oficial}\\[4pt]
\begin{tabular}{|c|c|c|c|c|c|c|}
\hline
1-2 & 3-4 & 5-6 & 7 & 8 & \textbf{TOTAL} & \textbf{NOTA} \\ \hline
 & & & & & & \\ \hline
\end{tabular}
\end{minipage}

\vspace{6pt}
\noindent\textbf{Instrucciones:} La duración del examen es de 1 hora y 50 minutos. El examen consta de ocho preguntas en tres hojas impresas por ambos lados, verifique que su examen esté completo. En las preguntas con procedimiento justifique sus respuestas en los espacios asignados. No está permitido sacar hojas en blanco ni ningún tipo de apuntes durante el examen, verifique que su celular esté apagado. No se permite el uso de calculadora.

\vspace{6pt}
\noindent\textbf{IDENTIFICACIÓN}

\vspace{8pt}
\noindent Nombre: \dotfill\ Cédula \dotfill

\vspace{4pt}
\noindent Profesor: \dotfill\ Grupo \dotfill

\vspace{6pt}
\begin{center}
\textbf{Espacio en blanco para realizar cómputos}
\end{center}

\newpage
\noindent\textbf{\large I. Completación}

\vspace{2pt}
\noindent En las preguntas 1 a 6 complete los espacios en blanco. \textbf{NOTA}: En esta sección se califica sólo la respuesta y no se tiene en cuenta el procedimiento.

\begin{enumerate}
\item{}[10pt] Sea $T:\mathbb{R}^3\to\mathbb{R}^3$ la transformación lineal dada por $T\left(\begin{bmatrix}x\\y\\z\end{bmatrix}\right)=\begin{bmatrix}x-y+z\\x-z\\2y+3z\end{bmatrix}$. Encontrar $[T]$, es decir, encontrar la matriz estándar de $T$.

\vspace{4pt}
\fbox{\begin{minipage}[t][3cm]{\dimexpr\linewidth-2\fboxsep-2\fboxrule}
$[T] = $
\end{minipage}}

\item{}[15pt] Determine si las siguientes afirmaciones son verdaderas o falsas (Marque V o F según el caso).

\vspace{6pt}
\noindent (a) [3pt] Si $S:\mathbb{R}^2\to\mathbb{R}^2$ es la transformación lineal dada por $S\left(\begin{bmatrix}x\\y\end{bmatrix}\right)=\begin{bmatrix}2y\\3x\end{bmatrix}$, entonces la inversa de $S$ está dada por $S^{-1}\left(\begin{bmatrix}x\\y\end{bmatrix}\right)=\begin{bmatrix}y/2\\x/3\end{bmatrix}$. \hfill\vf

\vspace{6pt}
\noindent (b) [3pt] El conjunto $W=\left\{\begin{bmatrix}x\\y\end{bmatrix}\in\mathbb{R}^2\ \middle|\ x=y\ \text{ó}\ x=-y\right\}$ es un subespacio de $\mathbb{R}^2$. \hfill\vf

\vspace{6pt}
\noindent (c) [3pt] El vector $\begin{bmatrix}1\\2\\-1\end{bmatrix}$ pertenece a $\mathrm{Col}\left(\begin{bmatrix}1&-1&1\\1&2&2\\0&1&1\end{bmatrix}\right)$. \hfill\vf

\vspace{6pt}
\noindent (d) [3pt] El vector $\begin{bmatrix}1\\2\\-1\end{bmatrix}$ pertenece a $\mathrm{Nul}\left(\begin{bmatrix}1&-1&1\\1&2&2\\0&1&1\end{bmatrix}\right)$. \hfill\vf

\vspace{6pt}
\noindent (e) [3pt] Si $p_1(x), p_2(x), p_3(x), p_4(x)$ son polinomios en el espacio vectorial $\mathcal{P}_3$, entonces el conjunto $\{p_1(x),p_2(x),p_3(x),p_4(x)\}$ es necesariamente L.D. \hfill\vf

\item{}[10pt] Sea $T:\mathcal{P}_2\to\mathcal{M}_{2\times 2}$ la transformación lineal dada por $T(a+bx+cx^2)=\begin{bmatrix}-a+b&a-c\\-b-c&c\end{bmatrix}$ y $S:\mathbb{R}^2\to\mathcal{P}_2$ la transformación lineal definida por $S\left(\begin{bmatrix}a\\b\end{bmatrix}\right)=(a+b)+bx+(2a-3b)x^2$. Calcular $T\circ S\left(\begin{bmatrix}2\\-1\end{bmatrix}\right)$.

\vspace{4pt}
\fbox{\begin{minipage}[t][3cm]{\dimexpr\linewidth-2\fboxsep-2\fboxrule}
$T\circ S\left(\begin{bmatrix}2\\-1\end{bmatrix}\right) = $
\end{minipage}}

\item{}[10pt] Sea $\mathcal{B}=\left\{\begin{bmatrix}1&1\\1&0\end{bmatrix},\begin{bmatrix}1&0\\1&1\end{bmatrix},\begin{bmatrix}0&1\\1&1\end{bmatrix}\right\}$ una base para un subespacio $Z$ del espacio vectorial $\mathcal{M}_{2\times 2}$. Encontrar una matriz $A$ tal que $[A]_{\mathcal{B}}=\begin{bmatrix}2\\0\\-1\end{bmatrix}$.

\vspace{4pt}
\fbox{\begin{minipage}[t][2.8cm]{\dimexpr\linewidth-2\fboxsep-2\fboxrule}
$A = $
\end{minipage}}

\item{}[10pt] Sea $\mathcal{B}=\{p_1(x)=x^3+2x,\ p_2(x)=x^2+2x+3,\ p_3(x)=x^3-2x+4,\ p_4(x)=x\}$. Determinar si $\mathcal{B}$ es una base para el espacio vectorial $\mathcal{P}_3$. Seleccione sí o no según el caso.

\vspace{6pt}
\begin{center}
\sinosino
\end{center}

\item{}[10pt] Sea $C = \begin{bmatrix}1&2&-1&0&1\\1&1&1&1&1\\2&3&0&1&2\end{bmatrix}$.
\begin{enumerate}
\item{}[5pt] Encuentre la dimensión de $\mathrm{Ren}(C)$.

\vspace{4pt}
\fbox{\begin{minipage}[t][2cm]{\dimexpr\linewidth-2\fboxsep-2\fboxrule}
$\dim(\mathrm{Ren}(C)) = $
\end{minipage}}

\item{}[5pt] Encuentre la nulidad de $C$, es decir, encuentre la dimensión del espacio nulo de $C$.

\vspace{4pt}
\fbox{\begin{minipage}[t][2cm]{\dimexpr\linewidth-2\fboxsep-2\fboxrule}
$\mathrm{Nulidad}(C) = $
\end{minipage}}
\end{enumerate}
\end{enumerate}

\vspace{6pt}
\noindent\textbf{\large II. Solución con Procedimiento}

\vspace{4pt}
\noindent En esta parte del examen se deben justificar las respuestas.

\begin{enumerate}
\setcounter{enumi}{6}
\item{}[10pt] Sea $A = \begin{bmatrix}1&1&0\\1&0&1\\1&0&0\end{bmatrix}$. Encontrar $A^{-1}$ en caso que $A$ sea invertible.

\vspace{6cm}

\newpage
\item{}[25pt] Determine una base y la dimensión para cada uno de los siguientes subespacios de $\mathbb{R}^4$.
\begin{enumerate}
\item{}[13pt] $V = \left\{\begin{bmatrix}-a+b\\a-d\\-b-c\\c+d\end{bmatrix}\in\mathbb{R}^4\ \middle|\ a,b,c,d\in\mathbb{R}\right\}$.

\vspace{6cm}

\item{}[12pt] $W = \left\{\begin{bmatrix}x\\y\\z\\w\end{bmatrix}\in\mathbb{R}^4\ \middle|\ -x+y=0,\ x-w=0,\ -y-z=0\right\}$.

\vspace{6cm}
\end{enumerate}
\end{enumerate}

\examfooter
\end{document}
